\chapter{Conclusion}

The present study supports the use of synthetic augmentation to improve model robustness in small-cohort Raman spectroscopy studies and highlights the potential of generative modeling as a complementary strategy for addressing data scarcity in biomedical spectral diagnostics. 

In this work, we addressed the problem of data scarcity in Raman spectroscopy-based glioma classification by developing a conditional variational autoencoder capable of generating class-conditioned synthetic spectra. We have designed and applied a spatio-spectral preprocessing pipeline to extract tumor-spectra from $58$ glioma biopsy slides, combining median filtering, NMF-based clustering, and GMM-guided visual selection. The conditional VAE was trained under a multi-task objective that jointly optimized spectral reconstruction and auxiliary class discrimination, with a KL-regularized latent space to support generative sampling. Synthetic spectra were evaluated under TS/TR and TSR/TR protocols within a strict slide-isolated cross-validation scheme across four downstream classifiers, i.e., RF+SVM, Random Forest, XGBoost, and $1$D-CNN, on both binary IDH-status and $6$-class methylation subtype tasks. We have also proposed and evaluated Classification by Reconstruction (CbR) as a generative inference strategy that bypasses discriminative classifiers entirely. We gained the insight that, whereas synthetic spectra were not able to substitute for real data, they are useful as an augmentation data tool in order to improve classification robustness.
